% Appendix E — Repository-selection pipeline (replication)

The 500 repositories used in this paper were selected by the procedure below. The script \texttt{scripts/01\_build\_repo\_list.py} implements it and writes \texttt{data/repos.json} and \texttt{results/phase1\_stats.json}. Executing the script fresh should reproduce the set up to GitHub's own indexing non-determinism (order of ties within a star-bucket and churn in repository metadata).

\paragraph{Goal.} Public GitHub repositories that plausibly gained $\geq 1{,}000$ stars during calendar year 2025 and were actively pushed within our analysis window, capped at the 500 most-active to keep Phase~2 PR collection tractable.

\paragraph{Step 1 — Star-bucket sweep (GitHub Search API).} The Search API returns at most 1{,}000 results per query. To cover the long tail we partition the star axis and query each bucket separately. Every query is
\begin{center}
\texttt{stars:LO..HI pushed:>=2025-04-01 is:public archived:false}
\end{center}
paginated \texttt{per\_page=100} up to 10 pages, sorted by \texttt{stars desc}. The 19 buckets used are:
\begin{center}\footnotesize
\begin{tabular}{@{}lrlr@{}}
\toprule
Bucket & \texttt{total\_count} seen & Bucket & \texttt{total\_count} seen \\
\midrule
1000--1099 & 1{,}000 & 4000--4999 & 1{,}000 \\
1100--1199 & 1{,}000 & 5000--6499 & 1{,}000 \\
1200--1299 & 1{,}000 & 6500--8999 & 1{,}000 \\
1300--1399 & 1{,}000 & 9000--14{,}999 & 1{,}000 \\
1400--1499 & 1{,}000 & 15{,}000--24{,}999 & 1{,}000 \\
1500--1699 & 1{,}000 & 25{,}000--49{,}999 & 864 \\
1700--1899 & 1{,}000 & 50{,}000--99{,}999 & 267 \\
1900--2099 & 1{,}000 & $\geq$100{,}000 & 94 \\
2100--2499 & 1{,}000 & & \\
2500--2999 & 1{,}000 & & \\
3000--3999 & 1{,}000 & & \\
\bottomrule
\end{tabular}
\end{center}
Buckets below the top three are saturated at 1{,}000 (GitHub reported total counts in the 1.5k--10k range and we capped at the API's hard 1{,}000-per-query limit). Total candidates collected across buckets, after deduplication by repo id: 17{,}225.

\paragraph{Step 2 — Non-fork / non-archived / non-mirror filter.} The Search query already excludes archived repositories, but we additionally drop \texttt{fork=true}. Forks and archives contribute noise without adding independent PR activity.

\paragraph{Step 3 — ``Gained $\geq$1{,}000 stars in 2025'' proxy.} GitHub's API does not directly expose historical star counts without expensive per-repo pagination of the stargazers endpoint. We use a conservative two-rule proxy:
\begin{enumerate}
\item \emph{Young-popular:} \texttt{created\_at} $\geq$ \texttt{2025-01-01} AND \texttt{stargazers\_count} $\geq$ \texttt{1{,}000}. Any repository satisfying this \emph{must} have gained $\geq 1{,}000$ stars in 2025, since it did not exist before.
\item \emph{Old-active-popular:} \texttt{created\_at} $<$ \texttt{2025-01-01} AND \texttt{stargazers\_count} $\geq$ \texttt{5{,}000} AND \texttt{pushed\_at} within the last 120 days at the time of sweep (2026-04-21). A repository with $\geq$5{,}000 stars and recent activity is very likely to have gained $\geq$1{,}000 stars in 2025 too, given typical open-source star-growth profiles. We choose this threshold strictly. It biases \emph{against} inclusion.
\end{enumerate}
Repositories are kept if \emph{either} rule fires. After this filter: 5{,}893 repositories.

\paragraph{Step 4 — Activity-score cap.} The 5{,}893 repositories are ranked by
\[
\mathrm{score}(r) = \log_{10}(\mathrm{stars}_r) \cdot (0.6 + 0.4\,\mathrm{recency}_r) + 0.3\,\mathbb{1}[\mathrm{created}_r \geq \mathtt{2025{-}01{-}01}],
\]
where $\mathrm{recency}_r = \exp(-\mathrm{days\_since\_push}_r / 60)$ is an exponential decay with $\approx$42-day half-life. We retain the top 500 by this score. Ties are broken by repository id (stable across re-runs).

\paragraph{Resulting sample characteristics.}
\begin{itemize}
\item 500 final repositories from 17{,}225 candidates $\to$ 5{,}893 after star filter $\to$ 500 after cap.
\item Keep-reason breakdown: 112 young-popular (created $\geq$ 2025-01-01), 388 old-active-popular.
\item Stars: min 18{,}755, median 53{,}583, max 457{,}447.
\item Top languages: TypeScript 120, Python 115, JavaScript 44, Go 40, Rust 39, C++ 19, Java 16, C 10, Shell 9, Ruby 9.
\item 173 GitHub API requests consumed across 19 Search-API buckets.
\end{itemize}

\paragraph{Known selection biases.} The activity-score cap \emph{concentrates the sample at the very-popular end} of GitHub: every repository we study has at least 18{,}755 stars. Our claims therefore apply to ``high-activity open-source'' and should not be generalised to GitHub at large. The ``$\geq$1{,}000 stars gained in 2025'' constraint is enforced only via the conservative proxy in Step~3. We expect a small fraction of old-active-popular repositories did \emph{not} actually gain that many stars in 2025. These repositories still satisfy the paper's implicit criterion (``popular-enough-to-have-AI-agent-activity in 2025--26'') and we therefore view the mismatch as harmless for the chain-length and DiD analyses.

\paragraph{Reproducing \texttt{data/repos.json}.} With a valid \texttt{GH\_TOKEN} (or \texttt{gh auth login}) and the subproject's \texttt{.venv} activated:
\begin{verbatim}
python scripts/01_build_repo_list.py \
    --cap 500 \
    --old-repo-star-threshold 5000 \
    --min-young-stars 1000
\end{verbatim}
Runtime: $\sim$13 minutes, dominated by Search API pacing (19 buckets $\times$ 10 pages $\times$ $\sim$1.2 s). Output: \texttt{data/repos.json} (one JSON object per repository, 500 total).

\paragraph{Sensitivity of the headline to the cap.} We do not claim the 3.4$\times$ year-on-year growth is invariant to the cap. A looser cap (e.g.\ 5{,}000 repos) would include more mid-popularity projects, likely lowering baseline AI-agent adoption in April 2025 and plausibly increasing the multiplier. A tighter cap (100 repos) would saturate on a handful of hyper-popular projects and lose statistical independence. The 500-repo choice was made to keep GraphQL budget feasible (Phase~2 cost $\approx$12{,}500 GraphQL requests over 7.4~hours) while preserving enough independent clusters for the cluster-bootstrap variance estimator.
